%----------
%	CONTENIDO NEGATIVO - MACHINE LEARNING
%----------	

We have followed a similar approach as in Section~\ref{sec:ageneral_ml}, with a total of 158 experiments in the span of five months (24/02/24 - 09/07/24), and defining five different windows illustrated in Figures ~\ref{fig:exp_cnegativo} and ~\ref{fig:exp_cnegativo_classes}. Figure \ref{fig:exp_cnegativo} shows the overall trend in weighted average F1 scores, and we can see that the global trend line, depicted in red, with an $R^2 = 0.37$, highlights a moderate yet positive impact of the changes made in each window. Figure \ref{fig:exp_cnegativo_classes}, on the other hand, breaks down the performance for specific classes: 'Insultos', 'Desprestigiar Acto', 'Desprestigiar Víctima', and 'Desprestigiar Deportista Autora'. This figure allows us to see how each class's F1 score evolves and reacts to different modifications. Each window also has its own trend line, providing a more granular view of performance over time.

\begin{figure}[H]
	\ffigbox[\FBwidth]
	{\caption{Distribution of the experiments conducted for "Contenido Negativo" over time with the different windows and trend line.}\label{fig:exp_cnegativo}}
 % [Aquí ponemos como queremos que aparezca en el índice, sin la referencia]{Aquí como queremos que aparezca en el documento, con la referencia}
	{\includegraphics[scale=0.25]{Plantilla_TFG_ingles_2019/imagenes/Experiments in _ContenidoNegativo_ over time.png}}
\end{figure}

\begin{figure}[H]
	\ffigbox[\FBwidth]
	{\caption{Distribution of the experiments conducted for "Contenido Negativo" per class over time with the different windows and trend line}\label{fig:exp_cnegativo_classes}}
 % [Aquí ponemos como queremos que aparezca en el índice, sin la referencia]{Aquí como queremos que aparezca en el documento, con la referencia}
	{\includegraphics[scale=0.25]{Plantilla_TFG_ingles_2019/imagenes/Experiments in _ContenidoNegativo_ over time_ perclass_ trendline.png}}
\end{figure}


\begin{table}[H]
    \centering
    \begingroup
    \small % Change to desired font size
    \ttabbox[\FBwidth]{
        \caption{Average Metrics for Each Window}
        \label{tab:window-metrics-cnegativo} % Unique label for the table
    }{
        \begin{tabularx}{\textwidth}{
            X
            >{\centering\arraybackslash}X
            >{\centering\arraybackslash}X
            >{\centering\arraybackslash}X
            >{\centering\arraybackslash}X
            >{\centering\arraybackslash}X
            >{\centering\arraybackslash}X
            >{\centering\arraybackslash}X
            >{\columncolor{gray!15}\centering\arraybackslash}X % Highlight this column
            >{\centering\arraybackslash}X
        }
            \hline
            \rowcolor{gray!30} % Adding a gray color to the header row
            \textbf{} & \textbf{Acc. (Std.)} & \textbf{Time (s)} & \textbf{W. Avg. Prec.} & \textbf{W. Avg. Rec.} & \textbf{W. Avg. F1-Score} & \textbf{D. Víc. F1-Score} & \textbf{D. Acto F1-Score} & \textbf{Ins. F1-Score} & \textbf{D. D. Aut. F1-Score} \\
            \hline
            W1 & 0.387 (0.053) & 289.778 & 0.399 & 0.387 & 0.369 & 0.376 & 0.343 & 0.316 & 0.417 \\
            \hline
            W2 & 0.512 (0.060) & 1262.331 & 0.512 & 0.512 & 0.506 & 0.472 & 0.410 & 0.358 & 0.639 \\
            \hline
            W3 & 0.458 (0.067) & 883.279 & 0.464 & 0.458 & 0.453 & 0.356 & 0.377 & 0.444 & 0.544 \\
            \hline
            W4 & 0.484 (0.051) & 757.250 & 0.493 & 0.484 & 0.465 & 0.337 & 0.347 & 0.381 & 0.618 \\
            \hline
            W5 & 0.576 (0.058) & 2680.062 & 0.592 & 0.576 & 0.579 & 0.641 & 0.405 & 0.465 & 0.685 \\
            \hline
            \multicolumn{10}{l}{Source: 'Training.xlsx'} \\
            \hline
        \end{tabularx}
    }
    \endgroup
\end{table}



\begin{itemize}
    \item \textbf{W1: Baseline Establishment (28-02-2024):} In this initial phase, the dataset remained unmodified, and traditional embeddings (Section~\ref{text_encoding_ml}) were used. Downsampling was used in order to test how the models behaved with data balancing techniques. While the best performance in terms of weighted F1-Score was 0.53, with all classes having F1 greater than 0.45, the average metrics for all classes had an F1-Score greater than 0.316, as seen in Table \ref{tab:window-metrics-cnegativo}.

    \item \textbf{W2: Introduction of Stylistic Features and Advanced Embeddings (08-04-2024):} During this period, several significant changes were introduced:
        \begin{itemize}
            \item \textbf{Data:} Stylistic features were added as described in Section~\ref{stylistic_features}, the 'Comentario Neutro' category was removed, and comprehensive text preprocessing was applied as detailed in Section~\ref{sec:data_preprocessing}.
            \item \textbf{Balancing:} SMOTE was employed to handle class imbalances.
            \item \textbf{Embedding:} New embeddings such as RoBERTa, BETO, and Robertuito were incorporated into the analysis.
        \end{itemize}
    From Table \ref{tab:window-metrics-cnegativo} we can see that he F1-scores during this window showed a varied response but consistent to these modifications, with improvements in F1-Score across all categories. In particular, we can see that the accuracy improved to 0.512 with a standard deviation of 0.060, but the average performance in terms of computation time was increased by 6.

    \item \textbf{W3: Comprehensive Testing of Modifications (26-04-2024):} This phase involved extensive testing of all prior modifications:
        \begin{itemize}
            \item \textbf{Balancing:} All balancing techniques were tested to identify the most effective method.
            \item \textbf{Embedding:} A comprehensive evaluation of all embeddings was conducted.
        \end{itemize}
    In this case, in Table \ref{tab:window-metrics-cnegativo} the performance metrics during this window revealed that F1-Score was not improved for all classes, unless for "Insultos". This is an important advancement as for the next classification task we are directly using tweets classified with this label.

    \item \textbf{W4: Control Phase Without Balancing Techniques (02-05-2024):} This window served as a control phase where all modifications were tested without applying any balancing techniques. From Figure \ref{fig:exp_cnegativo_classes} we can see that trend lines have a negative slope in this window, hence, resulting in a noticeable decline in model performance, as the only class that was improved was "Desprestigiar Deportista Autora" with F1-Score at 0.618.

    \item \textbf{W5: Introduction of OpenAI Embeddings (04-07-2024):} In the final phase, OpenAI embeddings were introduced, leading to the highest F1-scores observed throughout the experimentation in all classes. The overall weighted F1-Score reached 0.579, with "Desprestigiar Deportista Autora" achieving the highest F1-Score of 0.685, and 0.465 for "Insultos".


\end{itemize}

From this analysis, we provide the following tables with the best overall models for \textit{Contenido Negativo}. Table \ref{tab:cnegativo-bestmodels} lists the models and their correspondent parameters. Meanwhile, Table \ref{tab:cnegativo-bestmodels_metrics} shows the performance metrics for these models. Again, notice that for models {\cnegmajor} and {\cnegrandom} we are assuming that the Standard Deviation is zero because we are using the probabilities of each class, and that F1-Score is approximately equal to such probabilities for {\cnegrandom}.

\begin{table}[H]
    \centering
    \begingroup
    \small % Change to desired font size
    \ttabbox[\FBwidth]{
        \caption{Best Models for \textit{Contenido Negativo}}
        \label{tab:cnegativo-bestmodels} % Unique label for the table
    }{
        \begin{tabularx}{\textwidth}{
            X  % Name
            X  % Model
            X  % Embedding
            >{\hsize=0.5\hsize\centering\arraybackslas}X  % Embed. Size
            >{\hsize=0.3\hsize\centering\arraybackslas}X  % CV
            X  % Balance
            }
            \hline
            \rowcolor{gray!30} % Adding a gray color to the header row
            \textbf{} & \textbf{Model} & \textbf{Embedding} & \textbf{Embed. Size} & \textbf{CV} & \textbf{Balance} \\
            \hline
            \multicolumn{6}{l}{\textbf{Näive Models}} \\
            \hline
            \cnegmajor & - & - & - & - & - \\
            \hline
            \cnegrandom & - & - & - & - & - \\
            \hline
            \multicolumn{6}{l}{\textbf{Finetuned Models}} \\
            \hline
            \cneglrted & Logistic Regression & text-embedding-3-large & - & 5 & Down-sampling \\
            \hline
            \cnegrfbmd & Random Forest & BERT-multi & 500 & 5 & Down-sampling \\
            \hline
            \cnegxgbted & XGBoost & text-embedding-3-large & - & 5 & Down-sampling \\
            \hline
            \cnegxgbbua & XGBoost & Beto-uncased & - & 5 & ADASYN \\
            \hline
            \multicolumn{6}{l}{Source: 'Training.xlsx'} \\
            \hline
        \end{tabularx}
    }
    \endgroup
\end{table}


\begin{table}[H]
    \centering
    \begingroup
    \small % Change to desired font size
    \ttabbox[\FBwidth]{
        \caption{Metrics for the Best Models for \textit{Contenido Negativo}}
        \label{tab:cnegativo-bestmodels_metrics} % Unique label for the table
    }{
        \begin{tabularx}{\textwidth}{
            X
            >{\centering\arraybackslash}X
            >{\columncolor{gray!15}\centering\arraybackslash}X % Highlight this column
            >{\centering\arraybackslash}X
            >{\centering\arraybackslash}X
            >{\centering\arraybackslash}X
            >{\centering\arraybackslash}X
        }
            \hline
            \rowcolor{gray!30} % Adding a gray color to the header row
            \textbf{} & \textbf{Acc. (Std.)} & \textbf{Weighted Avg. F1-Score} & \textbf{D. Víc. F1-Score} & \textbf{D. Acto F1-Score} & \textbf{Ins. F1-Score} & \textbf{D. D. Aut. F1-Score} \\
            \hline
            \multicolumn{7}{l}{\textbf{Näive Models}} \\
            \hline
            \cnegmajor & 0.420 (0.000) & 0.420 & 0.000 & 0.000 & 0.000 & 1.000 \\
            \hline
            \cnegrandom & 0.264 (0.000) & 0.264 & 0.202 & 0.207 & 0.170 & 0.420 \\
            \hline
            \multicolumn{7}{l}{\textbf{Finetuned Models}} \\
            \hline
            \cneglrted & 0.750 (0.076) & 0.745 & 0.762 & 0.667 & 0.846 & 0.720 \\
            \hline
            \cnegrfbmd & 0.542 (0.129) & 0.519 & 0.210 & 0.545 & 0.695 & 0.625 \\
            \hline
            \cnegxgbted & 0.667 (0.113) & 0.666 & 0.643 & 0.632 & 0.692 & 0.696 \\
            \hline
            \cnegxgbbua & 0.479 (0.081) & 0.478 & 0.333 & 0.303 & 0.667 & 0.557 \\
            \hline
            \multicolumn{7}{l}{Source: 'Training.xlsx'} \\
            \hline
        \end{tabularx}
    }
    \endgroup
\end{table}


From this analysis we have decided to choose model \textbf{\cneglrted}  as the best performing model for this task, as it achieved the highest Weighted Average F1-Score of 0.750 while maintaining high performance across other classes with F1-Scores of 0.762, 0.667, 0.846 and 0.720 for "Desprestigiar Víctima", "Desprestigiar Acto", "Insultos" and "Desprestigiar Deportista Autora", respectively. It was really important to ensure a high Insultos F1-Score, as the third and final classification task, \textit{Insultos}, focuses exclusively on identifying negative tweets with insults. Moreover, it achieves the lowest standard deviation of the models detailed in Table~\ref{tab:cnegativo-bestmodels_metrics}, and significant improvements compared to näive models {\cnegmajor} and {\cnegrandom}.

We provide the confussion matrix for this model in Figure \ref{fig:exp_cnegativo_metrics}. In particular, we can see that the majority of errors in the classification occurred for "Desprestigiar Deportista Autora", "Desprestigiar Víctima" and "Desprestigiar Acto". One possible reason behind this is that the context and meaning of such classes do not differ too much between them, hence the misclassification. Moreover, for the class of interest, "Insultos", we have high precision of 0.786 and recall of 0.917.

\begin{figure}[H]
	\ffigbox[\FBwidth]
	{\caption{Confusion Matrix for the best overall model (\textbf{\cneglrted}) of the Machine Learning experiments conducted for "Contenido Negativo"}\label{fig:exp_cnegativo_metrics}}
 % [Aquí ponemos como queremos que aparezca en el índice, sin la referencia]{Aquí como queremos que aparezca en el documento, con la referencia}
	{\includegraphics[scale=0.65]{Plantilla_TFG_ingles_2019/imagenes/Experiments in _ContenidoNegativo_ best_model.png}}
\end{figure}